\documentclass[10pt,a4paper]{article}
\usepackage[utf8]{inputenc}
\usepackage[T1]{fontenc}
\usepackage[english,spanish]{babel}
\usepackage{plex-serif}\usepackage[sfdefault]{plex-sans}
\usepackage[margin=1.4cm]{geometry}
\usepackage{xcolor}\definecolor{navy}{HTML}{003A70}
\usepackage[colorlinks=true,urlcolor=navy]{hyperref}
\usepackage{titlesec}\usepackage{enumitem}
\titleformat{\section}{\serif\large\bfseries\color{navy}}{}{0pt}{}[\vspace{-0.6em}\color{navy}\rule{\linewidth}{0.8pt}]
\titlespacing*{\section}{0pt}{8pt}{4pt}
\setlist{nosep,leftmargin=1.2em}\setlength{\parindent}{0pt}\pagestyle{empty}
\newcommand{\serif}{\fontfamily{IBMPlexSerif-TLF}\selectfont}
\begin{document}
{\serif\Huge\bfseries\color{navy} Hideki Harada Oyakawa}\\[2pt]
{\small\color{gray} Politólogo · Educación y políticas públicas · Lima, Perú}\\
{\small \href{mailto:harada.oyakawa.hideki@gmail.com}{harada.oyakawa.hideki@gmail.com} $\cdot$ \href{https://haradaoyakawahideki.github.io}{haradaoyakawahideki.github.io} $\cdot$ \href{https://www.linkedin.com/in/hideki-harada}{linkedin.com/in/hideki-harada} $\cdot$ \href{https://github.com/haradaoyakawahideki}{github.com/haradaoyakawahideki}}
\section*{Educación}
\begin{itemize}
\item \textbf{PUCP, Q-Lab, Facultad de Ciencias Sociales.} Diplomatura de Especialización en Ciencia de Datos para las Ciencias Sociales y la Gestión Pública, ago. 2026 – may. 2027 (en curso)
\item \textbf{Davidson College}, Carolina del Norte, EE. UU. B.A. en Ciencia Política y Estudios de Género, 2021–2025 (grado conferido en mayo de 2026). \textit{Magna cum laude}, Altos Honores en Ciencia Política, becario UWC Davis.
\begin{itemize}
\item Intercambios: IES París (otoño 2023) · Semester at Sea, 10 países (primavera 2024)
\end{itemize}
\item \textbf{UWC Atlantic}, Gales, Reino Unido. Bachillerato Internacional, 2018–2020, becario UWC Davis.
\end{itemize}
\section*{Experiencia}
\smallskip
\textbf{Junior College Advisor}, Beca Cometa (iniciativa de Intercorp) · Lima · jul. 2026 – actualidad\par
\begin{itemize}
\item Asesoro a estudiantes peruanos de alto rendimiento y bajos recursos en su postulación a universidades de EE. UU. como estudiantes internacionales.
\item Acompañé a 55 estudiantes en sus solicitudes de ayuda financiera. Asesoro personalmente a 3 estudiantes durante todo su ciclo de postulación.
\end{itemize}
\smallskip
\textbf{Oficial Nacional de Selección}, Comité Nacional de UWC Perú · Remoto · oct. 2025 – actualidad\par
\begin{itemize}
\item Lidero y mejoro continuamente las cinco etapas del proceso nacional de selección de becarios.
\item Convocatoria 2027: más de 1,530 postulaciones de todas las regiones del Perú.
\item Convocatoria 2026: más de 600 postulaciones y más de 60 evaluadores voluntarios reclutados, capacitados y supervisados. El proceso resultó en 6 becas, 5 de ellas con financiamiento completo.
\end{itemize}
\smallskip
\textbf{Entrenador y Juez de Debate}, Liga Peruana de Debate Escolar (LPDE) · 2026\par
\begin{itemize}
\item Entrené al equipo ``Parque Nacional Huascarán'', \textbf{campeón mundial} del Mundial Escolar de Debate en Español (MED) 2026 en Panamá. El Perú ganó su cuarto título en una final íntegramente peruana. (\href{https://andina.pe/agencia/noticia-peru-se-corona-tetracampeon-el-mundial-debate-escolar-espanol-2026-1091241.aspx}{Andina})
\item Juez voluntario en el IV Torneo Descentralizado de Debate Puno 2026 en Juliaca (oct. 2026).
\end{itemize}
\smallskip
\textbf{Asistente de Investigación}, Departamento de Ciencia Política, Davidson College · may. – dic. 2025\par
\smallskip
\textbf{Pasante}, Oficina de Estudiantes Internacionales, Davidson College · abr. – ago. 2025\par
\begin{itemize}
\item Apoyé la orientación de más de 70 estudiantes internacionales y capacité a más de 20 líderes de orientación.
\end{itemize}
\smallskip
\textbf{Consultor}, Student Initiative for Academic Diversity, Davidson College · 2022 – 2025\par
\begin{itemize}
\item Entrevisté a candidatos docentes de cinco departamentos sobre sus prácticas pedagógicas inclusivas.
\end{itemize}
\smallskip
\textbf{Profesor Asistente y Tutor}, Estudios Hispánicos y Francófonos, Davidson College · 2022 – 2025\par
\smallskip
\textbf{Asistente de Investigación de Datos y DEI}, Centro de Desarrollo Vocacional, Davidson College · 2022 – 2023\par
\smallskip
\textbf{Tesorero}, International Student Association, Davidson College · 2022 – 2023\par
\begin{itemize}
\item Administré un presupuesto de más de US\$ 12,000 con 98 \% de ejecución. Las proyecciones ayudaron a conseguir 60 \% más de financiamiento.
\end{itemize}
\smallskip
\textbf{Pasante}, Association for Diplomatic Studies and Training · Washington DC · verano 2023\par
\smallskip
\textbf{Analista Junior}, CanopyLAB (EdTech) · Lima · 2020 – 2021\par
\section*{Reconocimientos}
\begin{itemize}
\item Equipo campeón mundial (como entrenador), Mundial Escolar de Debate en Español (MED), Panamá (2026)
\item Premio J. Harris Proctor en Ciencia Política, Davidson College (mayo 2026). Se otorga cada año desde 1993 a los estudiantes más destacados de la promoción en Ciencia Política.
\item Altos Honores en Ciencia Política, Davidson College (mayo 2026)
\item Fondo de investigación Davidson Research Initiative, US\$ 4,000
\item Primer puesto (equipo), Impactathon 2024, hackatón de políticas públicas de Impacta y la PUCP
\item Torneo Nacional de Debate Escolar, Perú (2015, 2016)
\end{itemize}
\section*{Idiomas y herramientas}
\begin{itemize}
\item \textbf{Idiomas:} español (nativo) · inglés (avanzado) · francés (conversacional) · lengua de señas americana (básico)
\item \textbf{Datos:} Python (análisis de encuestas, mapas, dashboards) · Stata y R (básico) · microdatos de encuestas (ENAHO, ENADES)
\item \textbf{Herramientas:} Claude Code · HubSpot · Salesforce · Asana · Notion
\end{itemize}
\end{document}
